\documentclass[11pt,a4paper]{article}
\usepackage{local-overleaf}

\hypersetup{
  pdftitle={Anima - Plan de privacidad y seguridad},
  pdfsubject={Gobernanza de datos, alineamiento HIPAA y ruta ISO/IEC 27001},
  pdfauthor={Equipo Anima}
}

\definecolor{Aqua}{HTML}{32B9D8}
\definecolor{DeepNavy}{HTML}{12364C}

\newcommand{\flowermark}{%
  \begin{tikzpicture}[baseline=-.7mm,x=1mm,y=1mm,scale=.55]
    \foreach \a in {0,60,...,300}{%
      \fill[Terracotta,rotate around={\a:(0,0)}] (0,2.7) ellipse[x radius=1.7,y radius=2.8];
    }
    \fill[Navy] (0,0) circle (1.15);
  \end{tikzpicture}%
}

\fancyhead[L]{\flowermark\hspace{1.2mm}\scriptsize\bfseries\color{Navy}ANIMA\hspace{1.5mm}\color{Line}|\hspace{1.5mm}\color{Muted}PRIVACIDAD Y SEGURIDAD}
\fancyhead[R]{\scriptsize\color{Muted}Documento de definición conjunta}
\fancyfoot[L]{\scriptsize\color{Muted}Equipo Anima · Municipalidad de la Ciudad de Salta}
\fancyfoot[R]{\scriptsize\color{Muted}\thepage\ / \pageref*{LastPage}}

\newcommand{\frameworkicon}[2]{%
  \tikz[baseline=-.7ex]{\node[circle,fill=#2,text=White,font=\bfseries\scriptsize,minimum size=8.5mm,inner sep=0pt]{#1};}%
}

\newcommand{\frameworkcard}[4]{%
  \begin{minipage}[t]{.315\textwidth}
    \begin{tcolorbox}[
      enhanced,colback=#1,colframe=#1,boxrule=0pt,arc=2mm,
      left=3mm,right=3mm,top=3mm,bottom=3mm,
      height=29mm,halign=center,valign=center
    ]
      {\small\bfseries\color{Navy}#3}\par\smallskip
      {\scriptsize\color{Navy}#4}
    \end{tcolorbox}
  \end{minipage}%
}

\newcommand{\principle}[3]{%
  \begin{minipage}[t]{.48\linewidth}
    \begin{softcard}[height=31mm,top=2.5mm,bottom=2.5mm]
      \begin{minipage}[t]{8mm}\vspace{0pt}\numberbadge{#1}\end{minipage}%
      \begin{minipage}[t]{\dimexpr\linewidth-9mm\relax}\vspace{0pt}
        {\footnotesize\bfseries\color{Navy}#2}\par\smallskip
        {\footnotesize #3}
      \end{minipage}
    \end{softcard}
  \end{minipage}%
}

\newcommand{\rolecard}[3]{%
  \begin{minipage}[t]{.315\textwidth}
    \begin{softcard}[height=43mm,top=2.5mm,bottom=2.5mm]
      {\scriptsize\bfseries\color{Terracotta}\MakeUppercase{#1}}\par
      {\small\bfseries\color{Navy}#2}\par\smallskip
      {\small #3}
    \end{softcard}
  \end{minipage}%
}

\newcommand{\workstream}[3]{%
  \begin{card}[height=38mm,top=2.5mm,bottom=2.5mm,before skip=3pt,after skip=3pt,fontupper=\footnotesize]
    \begin{minipage}[t]{8mm}\vspace{0pt}\numberbadge{#1}\end{minipage}%
    \begin{minipage}[t]{\dimexpr\linewidth-9mm\relax}\vspace{0pt}
      {\bfseries\color{Navy}#2}\par\smallskip
      #3
    \end{minipage}
  \end{card}%
}

\newcommand{\phasechip}[3]{%
  \begin{minipage}[t]{.188\textwidth}
    \begin{tcolorbox}[colback=#1!8,colframe=#1!35,boxrule=.5pt,arc=2mm,left=2mm,right=2mm,top=1.7mm,bottom=1.7mm,height=21mm,before upper=\centering]
      {\scriptsize\bfseries\color{#1}#2}\par
      {\scriptsize\color{Navy}#3}
    \end{tcolorbox}
  \end{minipage}%
}

\tikzset{
  dataflow/.style={rounded corners=2mm,draw=Line,fill=White,line width=.6pt,minimum width=36mm,minimum height=27mm,align=center,text=Navy,font=\small\bfseries},
  flowarrow/.style={-{Latex[length=2.2mm]},line width=.8pt,color=Muted}
}

\begin{document}
\RaggedRight

% 1
\kicker{07 · Condiciones de confianza}
\section{Privacidad y seguridad}
\sectionintro{Anima trabaja con información vinculada al bienestar emocional. El piloto debe proteger a cada persona y, al mismo tiempo, producir indicadores útiles para la gestión pública.}

\begin{navycard}[top=5mm,bottom=5mm]
  {\Large\bfseries La confianza se diseña antes de recibir datos reales.}\par\smallskip
  El Municipio debe aprobar la finalidad, los responsables, el circuito de consentimiento, los accesos y la respuesta ante incidentes. Sin esas definiciones, el piloto permanece en un entorno de demostración.
\end{navycard}

\vspace{4mm}
\noindent\makebox[\textwidth][c]{%
  \frameworkcard{Terracotta!14}{AR}{Base local}{Datos personales y de salud. Finalidad, consentimiento y seguridad.}%
  \hspace{4.65mm}%
  \frameworkcard{Sky}{H}{Alineamiento HIPAA}{Evaluación de riesgos y salvaguardas administrativas, físicas y técnicas.}%
  \hspace{4.65mm}%
  \frameworkcard{Sage!20}{ISO}{Ruta ISO 27001}{Gestión de seguridad, controles, evidencia y auditoría.}%
}%

\vspace{5mm}
\subsection{Seis compromisos que no se negocian}
\principle{1}{Finalidad definida}{Cada dato responde a una finalidad aprobada. No se reutiliza por conveniencia.}%
\hfill
\principle{2}{Consentimiento verificable}{La persona entiende qué se registra, para qué, por cuánto tiempo y cómo revocarlo.}

\principle{3}{Identidades separadas}{Los datos identificatorios se guardan fuera del repositorio analítico y con acceso restringido.}%
\hfill
\principle{4}{Lectura agregada}{El tablero municipal muestra grupos y zonas que superan un umbral mínimo. No expone casos.}

\principle{5}{Acceso trazable}{Cada consulta, cambio y exportación queda asociada a una persona, un rol y una fecha.}%
\hfill
\principle{6}{Respuesta humana}{La aplicación no diagnostica. Un equipo competente define y opera los protocolos de ayuda.}

\clearpage

% 2
\kicker{07.1 · Arquitectura de confianza}
\section{Cómo se protege el dato}
\sectionintro{La protección no depende de una sola base de datos. Depende de separar identidades, limitar los usos y evitar que el observatorio municipal pueda reconstruir historias individuales.}

\begin{center}
\begin{tikzpicture}[node distance=6mm]
  \node[dataflow,fill=Sand] (a) {1. Registro\\voluntario};
  \node[dataflow,fill=Terracotta!8,right=of a] (b) {2. Finalidad y\\consentimiento};
  \node[dataflow,fill=Sky!65,right=of b] (c) {3. Evento\\pseudonimizado};
  \node[dataflow,fill=Sage!17,right=of c] (d) {4. Agregación\\por umbral};
  \draw[flowarrow] (a)--(b);
  \draw[flowarrow] (b)--(c);
  \draw[flowarrow] (c)--(d);
  \node[dataflow,fill=Mist,below=8mm of b] (vault) {Bóveda de identidad\\separada y restringida};
  \node[dataflow,fill=Navy,text=White,below=8mm of d] (dash) {Observatorio municipal\\sin casos individuales};
  \draw[flowarrow] (b)--(vault);
  \draw[flowarrow] (d)--(dash);
  \draw[dashed,line width=.7pt,color=Terracotta] ($(vault.north east)+(2mm,2mm)$) -- ($(c.south west)+(-2mm,-2mm)$);
  \node[font=\scriptsize\bfseries,text=Terracotta,align=center] at ($(vault.north east)!0.5!(c.south west)+(0,-1mm)$) {separación\\lógica};
\end{tikzpicture}
\end{center}

\begin{navycard}[top=4mm,bottom=4mm]
  \whitecardtitle{Regla del observatorio}
  El tablero sólo recibe indicadores agregados. Si una combinación de zona, período y grupo no alcanza el umbral acordado, el sistema no muestra el resultado ni permite exportarlo.
\end{navycard}

\vspace{4mm}
\begin{minipage}[t]{.315\textwidth}
  \begin{card}[height=55mm,fontupper=\small]
    \cardtitle{Recolección}
    \begin{itemize}
      \item Sólo campos aprobados.
      \item Aviso de privacidad claro.
      \item Consentimiento con versión y fecha.
      \item Revocación y canal de consulta.
    \end{itemize}
  \end{card}
\end{minipage}\hfill
\begin{minipage}[t]{.315\textwidth}
  \begin{card}[height=55mm,fontupper=\small]
    \cardtitle{Custodia}
    \begin{itemize}
      \item Cifrado en tránsito y reposo.
      \item Acceso por rol y MFA.
      \item Copias y restauración probada.
      \item Retención y eliminación definidas.
    \end{itemize}
  \end{card}
\end{minipage}\hfill
\begin{minipage}[t]{.315\textwidth}
  \begin{card}[height=55mm,fontupper=\small]
    \cardtitle{Uso y control}
    \begin{itemize}
      \item Registro de accesos y exportaciones.
      \item Revisión periódica de permisos.
      \item Monitoreo de incidentes.
      \item Proveedores sujetos a contrato.
    \end{itemize}
  \end{card}
\end{minipage}

\vspace{5mm}
\begin{softcard}[height=39mm,top=3mm,bottom=3mm]
  \cardtitle{Pruebas antes de habilitar datos reales}
  \begin{minipage}[t]{.23\linewidth}\vspace{0pt}
    {\footnotesize\bfseries\color{Navy}Separación}\par
    {\footnotesize La vinculación con una identidad exige una función autorizada.}
  \end{minipage}\hfill
  \begin{minipage}[t]{.23\linewidth}\vspace{0pt}
    {\footnotesize\bfseries\color{Navy}Supresión}\par
    {\footnotesize Los grupos bajo el umbral no se muestran ni se exportan.}
  \end{minipage}\hfill
  \begin{minipage}[t]{.23\linewidth}\vspace{0pt}
    {\footnotesize\bfseries\color{Navy}Trazabilidad}\par
    {\footnotesize Cada acceso y exportación deja un registro auditable.}
  \end{minipage}\hfill
  \begin{minipage}[t]{.23\linewidth}\vspace{0pt}
    {\footnotesize\bfseries\color{Navy}Recuperación}\par
    {\footnotesize Una copia se restaura en una prueba aislada y documentada.}
  \end{minipage}
\end{softcard}

\clearpage

% 3
\kicker{07.2 · Gobernanza operativa}
\section{Quién decide y quién responde}
\sectionintro{La seguridad falla cuando una tarea queda sin dueño. El convenio del piloto debe asignar responsabilidades institucionales, técnicas y de cuidado antes de habilitar datos reales.}

\rolecard{Municipio}{Finalidad y decisión pública}{Aprueba el propósito, los perfiles habilitados, la retención y el protocolo institucional de respuesta.}%
\hfill
\rolecard{Equipo Anima}{Tecnología y operación}{Implementa los controles, conserva la evidencia técnica y comunica incidentes según el acuerdo.}%
\hfill
\rolecard{Equipo interdisciplinario}{Criterio de cuidado}{Define lenguaje, umbrales y rutas de ayuda. La plataforma no reemplaza su evaluación profesional.}

\vspace{5mm}
\begin{card}[top=3mm,bottom=3mm]
  \cardtitle{Matriz mínima de responsabilidades}
  \small
  \renewcommand{\arraystretch}{1.18}
  \begin{tabularx}{\linewidth}{@{}>{\bfseries\color{Navy}}p{43mm} >{\centering\arraybackslash}X >{\centering\arraybackslash}X >{\centering\arraybackslash}X@{}}
    \toprule
    \textcolor{Muted}{Decisión o tarea} & \textcolor{Muted}{Municipio} & \textcolor{Muted}{Anima} & \textcolor{Muted}{Equipo de cuidado} \\
    \midrule
    Finalidad y base de tratamiento & Aprueba & Documenta & Revisa \\
    Umbral de agregación & Aprueba & Implementa & Recomienda \\
    Perfiles y exportaciones & Aprueba & Configura & Consulta \\
    Retención y eliminación & Aprueba & Ejecuta & Consulta \\
    Incidente de seguridad & Decide comunicación & Contiene e informa & Evalúa impacto asistencial \\
    Ruta de ayuda & Habilita capacidad & Integra el canal & Define y opera \\
    \bottomrule
  \end{tabularx}
\end{card}

\vspace{4mm}
\subsection{Circuito de incidentes}
\begin{center}
\phasechip{Terracotta}{1 · DETECTAR}{Alerta y registro}%
\hfill
\phasechip{Terracotta}{2 · CONTENER}{Limitar el daño}%
\hfill
\phasechip{Sage}{3 · EVALUAR}{Datos y personas}%
\hfill
\phasechip{Aqua}{4 · COMUNICAR}{Autoridad y afectados}%
\hfill
\phasechip{Navy}{5 · CERRAR}{Corregir y aprender}
\end{center}

\vspace{2mm}
\begin{terracard}[top=3mm,bottom=3mm,fontupper=\footnotesize]
  \cardtitle{Cierre contractual antes del piloto}
  El convenio debe nombrar responsables, contactos de incidentes, retención, umbral de agregación, permisos de exportación y capacidad de respuesta humana para cada conjunto de datos.
\end{terracard}

\clearpage

% 4
\kicker{07.3 · Referencia internacional}
\section{Plan de alineamiento HIPAA}
\sectionintro{HIPAA sólo corresponde si interviene una entidad cubierta o un asociado de negocio alcanzado por la normativa estadounidense. Esa aplicabilidad se determina antes de hablar de cumplimiento.}

\begin{tcolorbox}[enhanced,colback=Aqua!8,colframe=Aqua!55,boxrule=.7pt,arc=2mm,left=4mm,right=4mm,top=3.5mm,bottom=3.5mm,borderline west={2.8pt}{0pt}{Aqua}]
  \frameworkicon{H}{Aqua}\hspace{2mm}{\bfseries\color{Navy}Estado propuesto: preparación y evidencia}\par\smallskip
  {\small Si HIPAA no resulta aplicable, el paquete de salvaguardas se conserva como referencia de seguridad para datos de salud.}
\end{tcolorbox}

\vspace{3mm}
\begin{minipage}[t]{.48\textwidth}
  \workstream{1}{Aplicabilidad y alcance}{Identificar entidades, flujos, datos protegidos, proveedores y contratos. Emitir una conclusión jurídica documentada.}
  \workstream{2}{Riesgo y responsabilidad}{Designar un responsable de seguridad. Inventariar activos y evaluar amenazas sobre confidencialidad, integridad y disponibilidad.}
  \workstream{3}{Salvaguardas administrativas}{Aprobar políticas, altas y bajas de acceso, capacitación, continuidad, revisión periódica y régimen de excepciones.}
\end{minipage}\hfill
\begin{minipage}[t]{.48\textwidth}
  \workstream{4}{Salvaguardas físicas y técnicas}{Proteger equipos y medios. Aplicar identidad única, MFA, mínimo privilegio, cifrado, auditoría, integridad y recuperación.}
  \workstream{5}{Proveedores y acuerdos}{Evaluar nube, mensajería y soporte. Incorporar obligaciones de seguridad y, si corresponde, acuerdos de asociado de negocio.}
  \workstream{6}{Incidentes y evidencia}{Definir detección, contención, evaluación, comunicación y cierre. Probar el circuito y conservar los registros.}
\end{minipage}

\vspace{4mm}
\begin{navycard}[top=4mm,bottom=4mm]
  \whitecardtitle{Criterios de salida}
  \begin{minipage}[t]{.48\linewidth}\vspace{0pt}
    \begin{itemize}[label=\textcolor{Aqua}{\rule{1.4mm}{1.4mm}}]
      \item Dictamen de aplicabilidad aprobado.
      \item Registro de riesgos con responsables.
      \item Matriz de controles y excepciones.
    \end{itemize}
  \end{minipage}\hfill
  \begin{minipage}[t]{.48\linewidth}\vspace{0pt}
    \begin{itemize}[label=\textcolor{Aqua}{\rule{1.4mm}{1.4mm}}]
      \item Contratos y políticas vigentes.
      \item Restauración de copias verificada.
      \item Simulacro de incidente documentado.
    \end{itemize}
  \end{minipage}
\end{navycard}

\begin{terracard}[top=3mm,bottom=3mm,fontupper=\small]
  \cardtitle{Declaración permitida en esta etapa}
  \textit{Anima adopta un plan de alineamiento basado en la HIPAA Security Rule.} No corresponde afirmar que el producto o el Municipio ya cumplen HIPAA.
\end{terracard}

\clearpage

% 5
\kicker{07.4 · Sistema de gestión}
\section{Ruta ISO/IEC 27001:2022}
\sectionintro{La norma exige un sistema de gestión de seguridad de la información. La certificación llega después de definir el alcance, operar los controles, reunir evidencia y superar una auditoría independiente.}

\begin{center}
\phasechip{Terracotta}{FASE 1 · SEM. 1-2}{Mandato y alcance}%
\hfill
\phasechip{Terracotta}{FASE 2 · SEM. 3-4}{Brechas y riesgos}%
\hfill
\phasechip{Sage}{FASE 3 · SEM. 5-10}{SGSI y controles}%
\hfill
\phasechip{Aqua}{FASE 4 · SEM. 11-12}{Auditoría interna}%
\hfill
\phasechip{Navy}{FASE 5 · POSTERIOR}{Certificación}
\end{center}

{\scriptsize\color{Muted}Cronograma de referencia. Se confirma cuando estén definidos el alcance, los proveedores y las integraciones del piloto.\par}

\vspace{4mm}
\begin{card}[top=3mm,bottom=3mm]
  \cardtitle{Trabajo y evidencia por fase}
  \small
  \renewcommand{\arraystretch}{1.10}
  \begin{tabularx}{\linewidth}{@{}>{\bfseries\color{Terracotta}}p{30mm} >{\RaggedRight\arraybackslash}X >{\RaggedRight\arraybackslash}p{47mm}@{}}
    \toprule
    \textcolor{Muted}{Fase} & \textcolor{Muted}{Trabajo} & \textcolor{Muted}{Evidencia de cierre} \\
    \midrule
    1. Alcance & Confirmar objetivos, roles, sedes, sistemas, datos y terceros. & Alcance del SGSI y responsables aprobados. \\
    2. Riesgos & Definir el método, valorar riesgos y priorizar tratamientos. & Informe de brechas y registro de riesgos. \\
    3. Implementación & Aprobar políticas, seleccionar controles y operar los procedimientos. & Declaración de aplicabilidad y registros de operación. \\
    4. Verificación & Medir, auditar internamente, corregir y revisar con la dirección. & Informe de auditoría, acciones y acta de revisión. \\
    5. Certificación & Contratar un organismo acreditado y atender sus hallazgos. & Resultado de auditoría y certificado, si corresponde. \\
    \bottomrule
  \end{tabularx}
\end{card}

\vspace{3mm}
\begin{minipage}[t]{.48\textwidth}
    \begin{softcard}[height=47mm,top=2.5mm,bottom=2.5mm,fontupper=\footnotesize]
    \cardtitle{Paquete documental}
    \begin{itemize}
      \item Política de seguridad y alcance.
      \item Método y registro de riesgos.
      \item Declaración de aplicabilidad.
      \item Gestión de accesos y proveedores.
      \item Continuidad y respuesta a incidentes.
      \item Auditoría y revisión de dirección.
    \end{itemize}
  \end{softcard}
\end{minipage}\hfill
\begin{minipage}[t]{.48\textwidth}
    \begin{softcard}[height=47mm,top=2.5mm,bottom=2.5mm,fontupper=\footnotesize]
    \cardtitle{Puerta de ingreso al piloto}
    \begin{itemize}
      \item Datos y finalidades aprobados.
      \item Riesgos altos tratados o aceptados.
      \item MFA, permisos y registros probados.
      \item Restauración de copias verificada.
      \item Protocolo de incidentes ensayado.
      \item Responsable de cada control designado.
    \end{itemize}
  \end{softcard}
\end{minipage}

\vspace{3mm}
\begin{navycard}[top=4mm,bottom=4mm]
  \whitecardtitle{Cómo se convierte esto en un presupuesto defendible}
  Se cotizan diagnóstico, arquitectura, controles, documentación, pruebas y auditoría como entregables verificables. La seguridad no debe aparecer como una suma genérica sin alcance.
\end{navycard}

\vspace{2mm}
{\scriptsize\color{Muted}
Fuentes oficiales: \href{https://www.argentina.gob.ar/normativa/nacional/ley-25326-2000-64790}{Ley 25.326};
\href{https://www.argentina.gob.ar/normativa/nacional/ley-26657-175977/texto}{Ley 26.657};
\href{https://www.argentina.gob.ar/normativa/nacional/resoluci%C3%B3n-47-2018-312662/texto}{Resolución AAIP 47/2018};
\href{https://www.hhs.gov/hipaa/for-professionals/security/laws-regulations/index.html}{HHS, HIPAA Security Rule};
\href{https://www.hhs.gov/hipaa/for-professionals/covered-entities/index.html}{HHS, aplicabilidad};
\href{https://www.iso.org/standard/27001}{ISO/IEC 27001:2022}. Referencia visual: \href{https://heyato.ai/}{Ato / Heyato}.\par}

\end{document}
